\section{Semaine 2 : Nettoyage et calcul de nouveaux déscripteurs}

Durant cette semaine, l'objectif était de débuter le traitement informatique des
données. Pour cela, j'ai utilisé le langage de programmation Python et
majoritairement la bibliothèque Pandas. Après avoir collecté les données sur
10 saisons et 5 championnats, la première étape a été de nettoyer les données
manquantes et de sélectionner les colonnes pertinentes du dataset.
J'ai notamment supprimer toutes les données relatives aux cotes des paris. En effet, 
ces colonnes sont supprimées pour éviter le Data-Leakage. 
Elles reflètent des probabilités de sortie qui fausseraient l'apprentissage réel du modèle.

Par la suite, j'ai construit de nouvelles caractéristiques à partir des données

brutes. En premier lieu, j'ai calculé le nombre de \textbf{jours de repos} séparant
le match à la date $t$ du match précédent à la date $t-1$. Ce descripteur peut
s'avérer peu pertinent pour les équipes ayant été reléguées puis promues, car
l'interruption entre deux saisons dans des championnats différents introduit un
biais dans le calcul.

J'ai également mis en place des indicateurs de \textbf{forme récente}, en calculant
pour chaque équipe la moyenne de buts marqués sur les 5 derniers matchs. De la même
manière, un indicateur de \textbf{force offensive} a été construit en moyennant le
nombre de tirs tentés sur les 5 derniers matchs, et un indicateur de
\textbf{force défensive} en moyennant le nombre de tirs subis sur la même fenêtre
glissante. Ces trois indicateurs ont été calculés aussi bien pour les statistiques
de fin de match (\textit{Full Time}) que pour les statistiques de mi-temps
(\textit{Half Time}).

Pour le \textbf{score Elo}\footnote{Implémentation inspirée de
\url{https://www.eloratings.net/about}}, chaque équipe part d'un score initial de
$1500$. Ce score est mis à jour après chaque match en fonction de l'issue du match
et de l'écart de buts, via un facteur multiplicatif $G$ qui amplifie la mise à jour
pour les victoires avec un grand écart. Deux scores Elo indépendants ont été
calculés : l'un basé sur les résultats finaux (\textit{FT\_Elo}), l'autre basé sur
les résultats à la mi-temps (\textit{HT\_Elo}), afin de capturer la capacité d'une
équipe à bien débuter les matchs.

Enfin, en s'inspirant de la démarche de Dixon \& Coles --- dont le modèle repose
sur le rapport $\alpha_i / \alpha_j$ entre les forces d'attaque --- et d'Egidi \&
Torelli --- qui construisent explicitement la différence de classement $w_n =
\text{rank}_H - \text{rank}_A$ comme prédicteur --- nous avons ajouté pour chaque
paire de descripteurs une colonne \textbf{\_diff} (différence $H - A$) et une
colonne \textbf{\_ratio} (rapport $H / A$). Ces nouvelles variables permettent au
modèle de disposer directement d'une mesure de l'écart de niveau entre les deux
équipes, sans avoir à l'inférer implicitement à partir des valeurs brutes.


\subsection*{Descriptif des features construites}

L'ensemble des caractéristiques construites est organisé en six groupes fonctionnels.

\begin{itemize}

    \item \textbf{Forme récente} --- \texttt{FT\_Hforme}, \texttt{FT\_Aforme},
    \texttt{HT\_Hforme}, \texttt{HT\_Aforme} : moyenne des buts marqués sur les
    5 derniers matchs, respectivement pour l'équipe à domicile (H) et l'équipe
    visiteuse (A). Cet indicateur est calculé aussi bien sur les statistiques de
    fin de match (\textit{Full Time}) qu'à la mi-temps (\textit{Half Time}).

    \item \textbf{Force offensive} --- \texttt{FT\_Hatt}, \texttt{FT\_Aatt} :
    moyenne du nombre de tirs tentés sur les 5 derniers matchs. Proxy de la
    pression offensive exercée par l'équipe.

    \item \textbf{Force défensive} --- \texttt{FT\_Hdef}, \texttt{FT\_Adef},
    \texttt{HT\_Hdef}, \texttt{HT\_Adef} : moyenne du nombre de tirs subis sur
    les 5 derniers matchs (\textit{Full Time}) et des buts concédés à la
    mi-temps (\textit{Half Time}). Plus cette valeur est faible, plus la défense
    de l'équipe est solide.

    \item \textbf{Précision} --- \texttt{FT\_Hprecision}, \texttt{FT\_Aprecision} :
    rapport entre le nombre moyen de tirs cadrés et le nombre moyen de tirs tentés
    sur les 5 derniers matchs. Mesure la qualité des occasions créées.
    La variante pondérée \texttt{FT\_Hprec\_weight}, \texttt{FT\_Aprec\_weight}
    donne un poids supplémentaire (facteur $w_g = 1{,}5$) aux tirs convertis en
    buts, afin de valoriser l'efficacité offensive.

    \item \textbf{Repos} --- \texttt{HRepos}, \texttt{ARepos} : nombre de jours
    écoulés depuis le dernier match joué par l'équipe (domicile ou extérieur),
    plafonné à 25 jours pour limiter les biais liés aux promotions et relégations.

    \item \textbf{Score Elo} --- \texttt{FT\_Elo\_H}, \texttt{FT\_Elo\_A},
    \texttt{HT\_Elo\_H}, \texttt{HT\_Elo\_A} : score Elo de chaque équipe avant
    le match, calculé itérativement dans l'ordre chronologique avec un score
    initial de $1500$. La mise à jour tient compte de l'issue et de l'écart de
    buts via le facteur $G$. Un score Elo indépendant est calculé sur les
    résultats à la mi-temps afin de capturer la tendance d'une équipe à bien ou
    mal démarrer les matchs.

    \item \textbf{Fautes et cartons} --- \texttt{FT\_HavgF}, \texttt{FT\_AavgF},
    \texttt{FT\_HavgY}, \texttt{FT\_AavgY}, \texttt{FT\_HavgR}, \texttt{FT\_AavgR} :
    moyennes respectivement du nombre de fautes commises, de cartons jaunes et de
    cartons rouges reçus sur les 5 derniers matchs. Ces indicateurs renseignent sur
    le style de jeu de l'équipe et sa discipline, les cartons rouges étant
    potentiellement corrélés aux matchs joués en infériorité numérique.

\end{itemize}

Pour chacun des groupes ci-dessus, trois colonnes dérivées sont systématiquement
ajoutées afin de faciliter la tâche des modèles de classification :

\begin{itemize}
    \item \textbf{\_diff} : différence $H - A$ entre la valeur de l'équipe à
    domicile et celle de l'équipe visiteuse. Par exemple,
    \texttt{FT\_forme\_diff} $=$ \texttt{FT\_Hforme} $-$ \texttt{FT\_Aforme}.
    Cette approche s'inspire directement d'Egidi \& Torelli, qui construisent
    $w_n = \text{rank}_H - \text{rank}_A$ comme prédicteur explicite de l'écart
    de niveau.

    \item \textbf{\_ratio} : rapport $H / A$ entre les deux valeurs. Cette
    formulation est cohérente avec la structure multiplicative du modèle de
    Dixon \& Coles, dans lequel la probabilité d'issue dépend du rapport
    $\alpha_i / \alpha_j$ entre les forces d'attaque des deux équipes.

    \item \textbf{\_dif} (pour l'Elo uniquement) : déjà inclus sous la forme
    \texttt{FT\_Elo\_dif} $=$ \texttt{FT\_Elo\_H} $-$ \texttt{FT\_Elo\_A} et
    \texttt{HT\_Elo\_dif}.
\end{itemize}

Au total, \textbf{53 features} sont ainsi construites et sauvegardées dans un fichier csv
, auxquelles s'ajoutent les trois colonnes cibles
\texttt{Hvs}, \texttt{Avs} et \texttt{Dvs} pour les analyses
\textit{One-vs-All}.

Trois colonnes cibles ont également été ajoutées --- \texttt{Hvs}, \texttt{Avs} et
\texttt{Dvs} --- afin de permettre une analyse selon le paradigme
\textit{One-vs-All} couramment utilisé en apprentissage automatique pour les
problèmes de classification multi-classes.